\section{AMC 10 10 2001}
        \begin{enumerate}
        	\item (\textit{AMC 10 10 2001}) The median of the list $n, n + 3, n + 4, n + 5, n + 6, n + 8, n + 10, n + 12, n + 15$
is $10$. What is the mean?


 $\textbf{(A) } 4 \qquad\textbf{(B) } 6 \qquad\textbf{(C) } 7 \qquad\textbf{(D) } 10 \qquad\textbf{(E) } 11$


 

	\item (\textit{AMC 10 10 2001}) A number $x$ is $2$ more than the product of its reciprocal and its additive inverse. In which interval does the number lie?


 $\textbf{(A) }  -4\leq x\leq -2 \qquad\textbf{(B) } -2<x\leq 0 \qquad\textbf{(C) } 0<x\leq 2 \qquad\textbf{(D) } 2<x\leq 4 \qquad\textbf{(E) } 4<x\leq 6$


 

	\item (\textit{AMC 10 10 2001}) The sum of two numbers is $S$. Suppose $3$ is added to each number and then each of the resulting numbers is doubled. What is the sum of the final two numbers?


 $\textbf{(A) } 2S+3 \qquad\textbf{(B) } 3S+2 \qquad\textbf{(C) } 3S+6 \qquad\textbf{(D) } 2S+6 \qquad\textbf{(E) } 2S+12$


 

	\item (\textit{AMC 10 10 2001}) What is the maximum number for the possible points of intersection of a circle and a triangle?


 $\textbf{(A) } 2 \qquad\textbf{(B) } 3 \qquad\textbf{(C) } 4 \qquad\textbf{(D) } 5 \qquad\textbf{(E) } 6$


 

	\item (\textit{AMC 10 10 2001}) How many of the twelve pentominoes pictured below have at least one line of reflectional symmetry?


 
                    \begin{figure}[H]
                        \centering
                        \includegraphics[width=0.5\textwidth]{Problem/AMC_10_10_2001_5_Problem_Figure_1.png}
                    \end{figure}
                    


 $\textbf{(A) } 3 \qquad\textbf{(B) } 4 \qquad\textbf{(C) } 5 \qquad\textbf{(D) } 6 \qquad\textbf{(E) } 7$


 

	\item (\textit{AMC 10 10 2001}) Let $P(n)$ and $S(n)$ denote the product and the sum, respectively, of the digits
of the integer $n$. For example, $P(23) = 6$ and $S(23) = 5$. Suppose $N$ is a
two-digit number such that $N = P(N)+S(N)$. What is the units digit of $N$?


 $\textbf{(A) } 2 \qquad\textbf{(B) } 3 \qquad\textbf{(C) } 6 \qquad\textbf{(D) } 8 \qquad\textbf{(E) } 9$


 

	\item (\textit{AMC 10 10 2001}) When the decimal point of a certain positive decimal number is moved four
places to the right, the new number is four times the reciprocal of the original
number. What is the original number?


 $\textbf{(A) } 0.0002 \qquad\textbf{(B) } 0.002 \qquad\textbf{(C) } 0.02 \qquad\textbf{(D) } 0.2 \qquad\textbf{(E) } 2$


 

	\item (\textit{AMC 10 10 2001}) Wanda, Darren, Beatrice, and Chi are tutors in the school math lab. Their
schedule is as follows: Darren works every third school day, Wanda works
every fourth school day, Beatrice works every sixth school day, and Chi works
every seventh school day. Today they are all working in the math lab. In how
many school days from today will they next be together tutoring in the lab?


 $\textbf{(A) } 42 \qquad\textbf{(B) } 84 \qquad\textbf{(C) } 126 \qquad\textbf{(D) } 178 \qquad\textbf{(E) } 252$


 

	\item (\textit{AMC 10 10 2001}) The state income tax where Kristin lives is levied at the rate of $p\%$ of the first $\textdollar 28000$ of annual income plus $(p + 2)\%$ of any amount above $\textdollar 28000$. Kristin noticed that the state income tax she paid amounted to $(p + 0.25)\%$ of her annual income. What was her annual income?


 $\textbf{(A) } \textdollar 28,000\qquad\textbf{(B) } \textdollar 32,000\qquad\textbf{(C) }\textdollar 35,000\qquad\textbf{(D) } \textdollar 42,000\qquad\textbf{(E) } \textdollar 56,000$


 

	\item (\textit{AMC 10 10 2001}) If $x$, $y$, and $z$ are positive with $xy = 24$, $xz = 48$, and $yz = 72$, then $x + y + z$ is


 $\textbf{(A) }18\qquad\textbf{(B) }19\qquad\textbf{(C) }20\qquad\textbf{(D) }22\qquad\textbf{(E) }24$


 

	\item (\textit{AMC 10 10 2001}) Consider the dark square in an array of unit squares, part of which is shown. The first ring of squares around this center square contains $8$ unit squares. The second ring contains $16$ unit squares. If we continue this process, the number of unit squares in the $100^\text{th}$ ring is


 
                    \begin{figure}[H]
                        \centering
                        \includegraphics[width=0.5\textwidth]{Problem/AMC_10_10_2001_11_Problem_Figure_1.png}
                    \end{figure}
                    


 $\textbf{(A)}\ 396 \qquad \textbf{(B)}\ 404 \qquad \textbf{(C)}\ 800 \qquad \textbf{(D)}\ 10,\!000 \qquad \textbf{(E)}\ 10,\!404$


 

	\item (\textit{AMC 10 10 2001}) Suppose that $n$ is the product of three consecutive integers and that $n$ is divisible by $7$. Which of the following is not necessarily a divisor of $n$?


 $\textbf{(A)}\ 6 \qquad \textbf{(B)}\ 14 \qquad \textbf{(C)}\ 21 \qquad \textbf{(D)}\ 28 \qquad \textbf{(E)}\ 42$


 

	\item (\textit{AMC 10 10 2001}) A telephone number has the form $ABC - DEF - GHIJ$, where each letter represents a different digit. The digits in each part of the numbers are in decreasing order; that is, $A > B > C$, $D > E > F$, and $G > H > I > J$. Furthermore, $D$, $E$, and $F$ are consecutive even digits; $G$, $H$, $I$, and $J$ are consecutive odd digits; and $A + B + C = 9$. Find $A$.


 $\textbf{(A)} \ 4 \qquad \textbf{(B)} \ 5 \qquad \textbf{(C)} \ 6 \qquad \textbf{(D)} \ 7 \qquad \textbf{(E)} \ 8$


 

	\item (\textit{AMC 10 10 2001}) A charity sells $140$ benefit tickets for a total of $\textdollar2001$. Some tickets sell for full price (a whole dollar amount), and the rest sells for half price. How much money is raised by the full-price tickets?


 $\textbf{(A) } \textdollar 782 \qquad \textbf{(B) } \textdollar 986 \qquad \textbf{(C) } \textdollar 1158 \qquad \textbf{(D) } \textdollar 1219 \qquad \textbf{(E) }\ \textdollar 1449$


 

	\item (\textit{AMC 10 10 2001}) A street has parallel curbs $40$ feet apart. A crosswalk bounded by two parallel stripes crosses the street at an angle. The length of the curb between the stripes is $15$ feet and each stripe is $50$ feet long. Find the distance, in feet, between the stripes.


 $\textbf{(A)}\ 9 \qquad \textbf{(B)}\ 10 \qquad \textbf{(C)}\ 12 \qquad \textbf{(D)}\ 15 \qquad \textbf{(E)}\ 25$


 

	\item (\textit{AMC 10 10 2001}) The mean of three numbers is $10$ more than the least of the numbers and $15$ less than the greatest. The median of the three numbers is $5$. What is their sum?


 $\textbf{(A)} \ 5 \qquad \textbf{(B)} \ 20 \qquad \textbf{(C)} \ 25 \qquad \textbf{(D)} \ 30 \qquad \textbf{(E)} \ 36$


 

	\item (\textit{AMC 10 10 2001}) Which of the cones listed below can be formed from a $252^\circ$ sector of a circle of radius $10$ by aligning the two straight sides?


 
                    \begin{figure}[H]
                        \centering
                        \includegraphics[width=0.5\textwidth]{Problem/AMC_10_10_2001_17_Problem_Figure_1.png}
                    \end{figure}
                    


 $\textbf{(A) } \text{A cone with slant height of 10 and radius 6}$


 $\textbf{(B) } \text{A cone with height of 10 and radius 6}$


 $\textbf{(C) } \text{A cone with slant height of 10 and radius 7}$


 $\textbf{(D) } \text{A cone with height of 10 and radius 7}$


 $\textbf{(E) } \text{A cone with slant height of 10 and radius 8}$


 

	\item (\textit{AMC 10 10 2001}) The plane is tiled by congruent squares and congruent pentagons as indicated. The percent of the plane that is enclosed by the pentagons is closest to


 
                    \begin{figure}[H]
                        \centering
                        \includegraphics[width=0.5\textwidth]{Problem/AMC_10_10_2001_18_Problem_Figure_1.png}
                    \end{figure}
                    


 $\textbf{(A)} \ 50 \qquad \textbf{(B)} \ 52 \qquad \textbf{(C)} \ 54 \qquad \textbf{(D)} \ 56 \qquad \textbf{(E)} \ 58 \qquad$


 

	\item (\textit{AMC 10 10 2001}) Pat wants to buy four donuts from an ample supply of three types of donuts: glazed, chocolate, and powdered. How many different selections are possible?


 $\textbf{(A)}\ 6 \qquad \textbf{(B)}\ 9 \qquad \textbf{(C)}\ 12 \qquad \textbf{(D)}\ 15 \qquad \textbf{(E)}\ 18$


 

	\item (\textit{AMC 10 10 2001}) A regular octagon is formed by cutting an isosceles right triangle from each of the corners of a square with sides of length $2000$. What is the length of each side of the octagon?


 $\textbf{(A) } \frac{1}{3}(2000) \qquad \textbf{(B) } {2000(\sqrt{2}-1)} \qquad \textbf{(C) } {2000(2-\sqrt{2})} \qquad \textbf{(D) } {1000} \qquad \textbf{(E) } {1000\sqrt{2}}$


 

	\item (\textit{AMC 10 10 2001}) A right circular cylinder with its diameter equal to its height is inscribed in a right circular cone. The cone has diameter $10$ and altitude $12$, and the axes of the cylinder and cone coincide. Find the radius of the cylinder.


 $\textbf{(A)}\ \frac{8}3\qquad\textbf{(B)}\ \frac{30}{11}\qquad\textbf{(C)}\ 3\qquad\textbf{(D)}\ \frac{25}{8}\qquad\textbf{(E)}\ \frac{7}{2}$


 

	\item (\textit{AMC 10 10 2001}) In the magic square shown, the sums of the numbers in each row, column, and diagonal are the same. Five of these numbers are represented by $v$, $w$, $x$, $y$, and $z$. Find $y + z$.


 
                    \begin{figure}[H]
                        \centering
                        \includegraphics[width=0.5\textwidth]{Problem/AMC_10_10_2001_22_Problem_Figure_1.png}
                    \end{figure}
                    


 $\textbf{(A)}\ 43 \qquad \textbf{(B)}\ 44 \qquad \textbf{(C)}\ 45 \qquad \textbf{(D)}\ 46 \qquad \textbf{(E)}\ 47$


 

	\item (\textit{AMC 10 10 2001}) A box contains exactly five chips, three red and two white. Chips are randomly removed one at a time without replacement until all the red chips are drawn or all the white chips are drawn. What is the probability that the last chip drawn is white?


 $\textbf{(A)}\ \frac{3}{10}\qquad\textbf{(B)}\ \frac{2}{5}\qquad\textbf{(C)}\ \frac{1}{2}\qquad\textbf{(D)}\ \frac{3}{5}\qquad\textbf{(E)}\ \frac{7}{10}$


 

	\item (\textit{AMC 10 10 2001}) In trapezoid $ABCD$, $\overline{AB}$ and $\overline{CD}$ are perpendicular to $\overline{AD}$, with $AB+CD=BC$, $AB<CD$, and $AD=7$. What is $AB\cdot CD$?


 $\textbf{(A)}\ 12 \qquad \textbf{(B)}\ 12.25 \qquad \textbf{(C)}\ 12.5 \qquad \textbf{(D)}\ 12.75 \qquad \textbf{(E)}\ 13$


 

	\item (\textit{AMC 10 10 2001}) How many positive integers not exceeding $2001$ are multiples of $3$ or $4$ but not $5$?


 $\textbf{(A)}\ 768 \qquad \textbf{(B)}\ 801 \qquad \textbf{(C)}\ 934 \qquad \textbf{(D)}\ 1067 \qquad \textbf{(E)}\ 1167$


 

\end{enumerate}